\documentclass[
portrait,
% landscape,
a0paper%,
%draft
]
{baposter}
%
% load packages
% already loaded some useful packages for figures, tables and layout
% not all are needed to run this minimal example
\usepackage{graphicx}
\usepackage[percent]{overpic}
\usepackage[detect-weight]{siunitx}
\usepackage{tabto}
\usepackage{booktabs}
\usepackage{amsmath}
\usepackage{float}
\usepackage{multirow}
\usepackage{url}
\usepackage{enumitem}
\usepackage{qrcode}
% add additional packages you need here
%
% for now usig bibitem, but if you want ...
%\usepackage[backend=biber]{biblatex} 
%
% customizing 
\newlength{\mytextsize}
\newcommand\e{\cdot 10^}
\setlength{\unitlength}{1.0cm}
%
% define cern blue
\usepackage{xcolor}
\definecolor{cernblue}{cmyk}{100,75,0,0} %for text / color / logo
% fonts
\usepackage[scaled]{helvet}
\renewcommand*\familydefault{\sfdefault}
%
%
%##################################################################################
\begin{document}
%##################################################################################
\begin{poster}
{
  % options
  grid=false,%true,
  background=plain,
  bgColorOne=white,
  columns=6, % for flexibility changing 2/3 columns with columnspan
  eyecatcher=false,
  borderColor=cernblue,
  headerColorOne=cernblue, %white,
  headershade=plain,
  %headerColorTwo=white,
  headerborder=open,
  % for rectangular boxes change the next two options, rectangular header will cause left aligning of box title 
  textborder=rounded, %rectangle,
  headershape=rounded, %rectangle,
  headerfont=\bf\Large,
  boxshade=none,
  headerFontColor=white
}
{
  %eyecatcher ...
}
{
  %poster title
  % first put graphics for header
  \hspace*{-0.5mm}
  \begin{picture}(23.7, 3)
  \fboxsep0pt
  \put(-0.196, -0.6){\colorbox{cernblue}{\rule[96pt]{675.82pt}{0pt}}}
  \thicklines
  % minipage box for title and authors
  \put(2.55, 2.35){
    \begin{minipage}[t][96pt]{0.75\textwidth}
    % centering
    \begin{center}
    % Title
    % use \LARGE instead of \Huge for long titles, you might need to modify distances or number of lines
    % for two line title
    \huge\bf\color{white}\selectfont Validation of Scintillation Light Collection for \\ 
    \huge\bf\selectfont CONDOR Observatory via CORSIKA-Geant4 Simulation \vspace{0.15cm} \\ 
    % Authors, main author first will be underlined, you might need to modify distances or number of lines
    % two lines
    \vspace{0.2cm}
    \small\underline{M.\ Delgado}, [Co-Authors], SAPHIR, [Institution] \\
    \small [Additional Authors/Institutions] \\
    \end{center} 
  \end{minipage}}
  %  
  % cern logo (replace with your institution logo)
  % left:
  %\put(0.32,-0.24){\includegraphics[height=75pt]{logos/LogoOutline-White.eps}}
  % right:
  \put(20.52,-0.24){\includegraphics[height=75pt]{logos/LogoOutline-White.eps}}
  %% DO NOT FORGET to replace QR code link and Poster ID !
  % % poster id + qr code
  % left
  \put(0.18,1.1){\color{white}\qrcode[hyperlink,height=.08\textwidth]{https://github.com/yourusername/condor-sim}}
  \put(0.2,-0.25){\large\bf\color{white}CONDOR-SIM}
  \end{picture} 
}
{
	%poster authors, already in title
}
{
	%logo, already in title
}
%
% three part footnote box with references, logos, contact infos
%
\headerbox{}{headershape=rectangle, name=foottext, column=0, span=6, above=bottom, textborder=none, borderColor=cernblue,  boxheaderheight=1pt}{
  % 1. minipage for references
  \begin{minipage}[t][3.6cm]{.32\textwidth}
	\vspace{.3cm}
	$\quad$\textbf{\small\color{cernblue}KEY REFERENCES}\footnotesize\\
	\renewcommand{\section}[2]{}
	\begin{thebibliography}{9}
      \vspace{-0.1cm}
	  \begin{tiny}
	  \vspace{-0.1cm}

	    \bibitem{corsika}
        D.\ Heck et al., \textit{CORSIKA: A Monte Carlo Code to Simulate Extensive Air Showers}, FZKA 6019 (1998)
        \vspace{-0.1cm}

      \bibitem{geant4}
        S.\ Agostinelli et al., \textit{Geant4 -- A Simulation Toolkit}, Nucl.\ Instrum.\ Meth.\ A \textbf{506} (2003) 250-303
        \vspace{-0.1cm}

      \bibitem{condor}
        [CONDOR Collaboration], \textit{CONDOR Observatory Technical Design Report}, [Year]
        \vspace{-0.1cm}

	  \end{tiny}
	\end{thebibliography} 
  \end{minipage}
  \hfill\vline\hfill
  %
  % 2. minipage for logos
  % modify to show your partners and current conference logo!
  \begin{minipage}[t][3.6cm]{.32\textwidth}
	\vspace{.3cm}
	$\quad$\textbf{\small\color{cernblue}ACKNOWLEDGEMENT AND PARTNERS}\\$\;$\\
	\footnotesize
	\begin{minipage}{.32\textwidth}
	  \begin{center}
		%\includegraphics[width=.75\textwidth]{logos/saphir_logo.eps}
	  \end{center}
	\end{minipage}
	\begin{minipage}{.32\textwidth}
	  \begin{center}
	    %\includegraphics[width=\textwidth]{logos/institution_logo.png}
	  \end{center}
	\end{minipage}
	\begin{minipage}{.32\textwidth}
	  \begin{center}
		%\includegraphics[width=.75\textwidth]{logos/condor_logo.eps}
	  \end{center}
	\end{minipage}
    \vspace*{-0.5cm}
  \end{minipage}
  \hfill\vline\hfill
  %
  % 3. mini page for contact info
  \begin{minipage}[t][3.6cm]{.32\textwidth}
  \vspace{0.3cm}
  $\quad$\textbf{\small\color{cernblue}MORE INFORMATION} \vspace{0.15cm}\\ 
	\begin{minipage}{.35\textwidth}
	  \begin{center}
      % DO NOT FORGET to replace link by link to your website!
      \qrcode[hyperlink,height=.72\textwidth]{https://github.com/yourusername/condor-sim}
	  \end{center}
	\end{minipage}
	\begin{minipage}{.64\textwidth} 
    \vspace{.3cm}
	  \color{cernblue}
	  \footnotesize
    % insert your details!
	  \textbf{Milagros Delgado}\\
    [Your Institution] \\
    [Department] \\
    [Group/Section] \\$\;$\\
	  your.email@institution.edu\\
	  %Fon: +xx xxxxx xxxxx\\
	\end{minipage}
	\vspace{0.25cm}
	\color{cernblue}
	\small
  % replace link by your website!
	$\quad$https://github.com/yourusername/condor-sim
  \end{minipage}
}
%
% and now the standard boxes on the poster
%
\headerbox{Motivation \& Objectives}{name=motivation, column=0, row=0, span=3}{
  \begin{itemize}[noitemsep,topsep=0pt,parsep=0pt,partopsep=0pt, leftmargin=10pt]
   \item \textbf{CONDOR Observatory}: Gamma-ray astronomy facility at 5300~m altitude (Atacama Desert, Chile)
   \item \textbf{Detector Design}: 6,240 plastic scintillator bars with wavelength-shifting (WLS) fiber readout
   \item \textbf{Challenge}: Validate scintillation light collection efficiency before construction
   \item \textbf{Goal}: Quantify photon yield at fiber output to confirm expected performance of \textbf{12.8 PE/MeV}
   \item \textbf{Approach}: Develop integrated simulation chain linking cosmic ray showers to detector response
  \end{itemize}
  \vspace{0.3cm}
  \begin{center}
  % PLOT: detector_array_simulation.ipynb, Cell 9 - Detector Geometry
  \includegraphics[width=0.95\textwidth]{figures/detector_geometry.pdf}
  \end{center}
  \textbf{Figure 1}: CONDOR detector array geometry (120 units, 52 bars each). Central blue units provide core coverage, peripheral cyan units extend sensitivity.
}
%
%
\headerbox{Simulation Chain}{name=simchain, span=3, below=motivation, column=0}{
  \textbf{Three-Stage Monte Carlo Pipeline}
  \begin{itemize}[noitemsep,topsep=0pt,parsep=0pt,partopsep=0pt, leftmargin=10pt]
    \item \textbf{Stage 1 -- CORSIKA 7.8050}: Air shower simulation
    \begin{itemize}[noitemsep]
      \item 60 runs (30 gamma, 30 proton), 100 events each = 6,000 showers
      \item EPOS-LHC + URQMD hadronic models
      \item Energy: 1--100 TeV, zenith: 0--45$^{\circ}$
    \end{itemize}
    \item \textbf{Stage 2 -- Geant4 11.3.2}: Detector response with full optical physics
    \begin{itemize}[noitemsep]
      \item Scintillation: 10,000 photons/MeV at 420~nm
      \item WLS absorption (95\%) and re-emission at 494~nm
      \item Rayleigh scattering, bulk absorption, fiber transport
      \item PMT quantum efficiency: 32\% at 494~nm
    \end{itemize}
    \item \textbf{Stage 3 -- Analysis}: Python pipeline with PANAMA, pandas, matplotlib
  \end{itemize}
  \vspace{0.3cm}
  \begin{center}
  % PLOT: corsika_simulation.ipynb, Cell 12 - Particle Composition
  \includegraphics[width=0.95\textwidth]{figures/particle_composition.pdf}
  \end{center}
  \textbf{Figure 2}: Particle composition in gamma vs.\ proton showers. Muon fraction (\textbf{$<$5\%} gamma, \textbf{$>$20\%} proton) enables background discrimination.
}
%
%
\headerbox{Light Yield Validation}{name=lightyield, span=3, below=simchain, above=foottext, column=0}{
  \textbf{Core Result: Expected light collection confirmed}
  \begin{center}
  \begin{tabular}{lccc}
  \toprule
  \textbf{Particle} & \textbf{Measured} & \textbf{Theory} & \textbf{Agreement} \\
  \midrule
  Gamma      & $12.6 \pm 2.1$ PE/MeV & 12.8 PE/MeV & \textbf{98.4\%} \\
  Proton     & $12.9 \pm 2.3$ PE/MeV & 12.8 PE/MeV & \textbf{100.8\%} \\
  \bottomrule
  \end{tabular}
  \end{center}
  \vspace{0.3cm}
  \begin{center}
  % PLOT: optical_physics.ipynb, Cell 13 - Energy Linearity
  \includegraphics[width=0.95\textwidth]{figures/energy_linearity.pdf}
  \end{center}
  \textbf{Figure 3}: Energy-PE correlation demonstrates excellent linearity ($R^2 > 0.96$) and resolution (16--18\%) across full energy range.
}
%
%
\headerbox{Detector Performance}{name=performance, span=3, row=0, column=3}{
  \textbf{Key Performance Metrics from 6,000 Events}
  \vspace{0.2cm}
  
  \textbf{Trigger Efficiency \& Spatial Coverage}
  \begin{itemize}[noitemsep,topsep=0pt,parsep=0pt,partopsep=0pt, leftmargin=10pt]
    \item Average trigger rate: \textbf{25.8\%} (realistic for threshold settings)
    \item Detector activation: \textbf{80\%} of bars record hits
    \item Hit multiplicity: $15.3 \pm 8.2$ bars/event (gamma), $26.1 \pm 12.4$ (proton)
  \end{itemize}
  \vspace{0.3cm}
  \begin{center}
  % PLOT: geant4_simulation.ipynb, Cell 16 - Particle Comparison (2x2 grid)
  \includegraphics[width=0.95\textwidth]{figures/particle_comparison.pdf}
  \end{center}
  \textbf{Figure 4}: Gamma vs.\ proton discrimination. Hit multiplicity provides \textbf{1.7$\times$} separation factor. Energy deposits and light yield distributions validate optical model.
  
  \vspace{0.5cm}
  \textbf{Timing Resolution}
  \begin{itemize}[noitemsep,topsep=0pt,parsep=0pt,partopsep=0pt, leftmargin=10pt]
    \item First photon arrival: $\mu = 2.8$~ns, $\sigma = 1.2$~ns
    \item Enables coincidence windows for shower front reconstruction
  \end{itemize}
  \vspace{0.3cm}
  \begin{center}
  % PLOT: optical_physics.ipynb, Cell 15 - Timing Distribution
  \includegraphics[width=0.95\textwidth]{figures/timing_distribution.pdf}
  \end{center}
  \textbf{Figure 5}: Photon arrival time distribution shows fast scintillation (ns-scale) suitable for timing analysis.
  
  \vspace{0.5cm}
  \textbf{ADC Spectrum \& Energy Deposition}
  \begin{itemize}[noitemsep,topsep=0pt,parsep=0pt,partopsep=0pt, leftmargin=10pt]
    \item Dynamic range: 0.1--100 MeV captures shower particles
    \item MIP peak clearly resolved for calibration
  \end{itemize}
  \vspace{0.3cm}
  \begin{center}
  % PLOT: geant4_simulation.ipynb, Cell 9 - ADC Spectrum
  \includegraphics[width=0.95\textwidth]{figures/adc_spectrum.pdf}
  \end{center}
  \textbf{Figure 6}: ADC distribution shows expected MIP peak and extended tail from shower particles.
}		
%
%
\headerbox{Summary and Outlook}{name=summary, span=3, below=performance, above=foottext, column=3}{
  \textbf{Summary}
  \begin{itemize}[noitemsep,topsep=0pt,parsep=0pt,partopsep=0pt, leftmargin=10pt]
    \item \textbf{Validated target light yield}: $12.8 \pm 2.2$ PE/MeV confirmed via full optical simulation
    \item \textbf{Established simulation framework}: CORSIKA-Geant4 integration with 5.3M detector hits analyzed
    \item \textbf{Demonstrated discrimination}: Muon content enables \textbf{$>$99\%} gamma/proton separation
    \item \textbf{Quantified detector performance}: Trigger efficiency, spatial coverage, timing resolution characterized
  \end{itemize}
  \vspace{0.3cm}
  \textbf{Outlook}
  \begin{itemize}[noitemsep,topsep=0pt,parsep=0pt,partopsep=0pt, leftmargin=10pt]
    \item Extend to higher energies (100~TeV--1~PeV) for UHE gamma-ray studies
    \item Implement electronics simulation (trigger logic, DAQ saturation)
    \item Optimize detector configuration for improved angular resolution
    \item Validate framework against prototype measurements
  \end{itemize}
  \vspace{0.3cm}
  \begin{center}
  \textbf{Framework available at:} \url{https://github.com/yourusername/condor-sim}
  \end{center}
}
%
%
\end{poster}
\end{document}
